% Options for packages loaded elsewhere
\PassOptionsToPackage{unicode}{hyperref}
\PassOptionsToPackage{hyphens}{url}
\PassOptionsToPackage{dvipsnames,svgnames,x11names}{xcolor}
\documentclass[
  11pt,
  a4paper,
]{article}
\usepackage{xcolor}
\usepackage[margin=2.6cm]{geometry}
\usepackage{amsmath,amssymb}
\setcounter{secnumdepth}{5}
\usepackage{iftex}
\ifPDFTeX
  \usepackage[T1]{fontenc}
  \usepackage[utf8]{inputenc}
  \usepackage{textcomp} % provide euro and other symbols
\else % if luatex or xetex
  \usepackage{unicode-math} % this also loads fontspec
  \defaultfontfeatures{Scale=MatchLowercase}
  \defaultfontfeatures[\rmfamily]{Ligatures=TeX,Scale=1}
\fi
\usepackage{lmodern}
\ifPDFTeX\else
  % xetex/luatex font selection
  \setmainfont[]{Libertinus Serif}
  \setmonofont[Scale=0.85]{Latin Modern Mono}
  \setmathfont[]{Latin Modern Math}
\fi
% Use upquote if available, for straight quotes in verbatim environments
\IfFileExists{upquote.sty}{\usepackage{upquote}}{}
\IfFileExists{microtype.sty}{% use microtype if available
  \usepackage[]{microtype}
  \UseMicrotypeSet[protrusion]{basicmath} % disable protrusion for tt fonts
}{}
\makeatletter
\@ifundefined{KOMAClassName}{% if non-KOMA class
  \IfFileExists{parskip.sty}{%
    \usepackage{parskip}
  }{% else
    \setlength{\parindent}{0pt}
    \setlength{\parskip}{6pt plus 2pt minus 1pt}}
}{% if KOMA class
  \KOMAoptions{parskip=half}}
\makeatother
\usepackage{longtable,booktabs,array}
\newcounter{none} % for unnumbered tables
\usepackage{calc} % for calculating minipage widths
% Correct order of tables after \paragraph or \subparagraph
\usepackage{etoolbox}
\makeatletter
\patchcmd\longtable{\par}{\if@noskipsec\mbox{}\fi\par}{}{}
\makeatother
% Allow footnotes in longtable head/foot
\IfFileExists{footnotehyper.sty}{\usepackage{footnotehyper}}{\usepackage{footnote}}
\makesavenoteenv{longtable}
\setlength{\emergencystretch}{3em} % prevent overfull lines
\providecommand{\tightlist}{%
  \setlength{\itemsep}{0pt}\setlength{\parskip}{0pt}}

\providecommand{\E}{\mathbb{E}}
\providecommand{\N}{\mathbb{N}}
\providecommand{\Z}{\mathbb{Z}}
\providecommand{\R}{\mathbb{R}}
\providecommand{\eps}{\varepsilon}
\providecommand{\ceil}[1]{\left\lceil #1\right\rceil}
\providecommand{\floor}[1]{\left\lfloor #1\right\rfloor}
\AtBeginDocument{\let\setminus\smallsetminus}
\usepackage[most]{tcolorbox}
\newtcolorbox{warning}{colback=red!5,colframe=red!65!black,boxrule=0.8pt,arc=2pt,left=6pt,right=6pt,top=5pt,bottom=5pt}
\usepackage{etoolbox}
\AtBeginEnvironment{longtable}{\small}
\setlength{\emergencystretch}{3em}
\usepackage{fancyhdr}
\pagestyle{fancy}
\fancyhf{}
\fancyhead[L]{\small\itshape Candidate result 1602.05184\_\_00 --- GPT-6 Astra, awaiting review}
\fancyhead[R]{\small\thepage}
\renewcommand{\headrulewidth}{0.2pt}
\usepackage{bookmark}
\IfFileExists{xurl.sty}{\usepackage{xurl}}{} % add URL line breaks if available
\urlstyle{same}
\hypersetup{
  pdftitle={A candidate proof of Conjecture 5 in On the difference between the Szeged and Wiener index},
  pdfauthor={github.com/graph-theory-AI --- machine-generated candidate writeup},
  colorlinks=true,
  linkcolor={blue!50!black},
  filecolor={Maroon},
  citecolor={Blue},
  urlcolor={blue!50!black},
  pdfcreator={LaTeX via pandoc}}

\title{A candidate proof of Conjecture 5 in On the difference between
the Szeged and Wiener index}
\author{\href{https://github.com/graph-theory-AI}{github.com/graph-theory-AI}
--- machine-generated candidate writeup}
\date{17 September 2026}

\begin{document}
\maketitle
\begin{abstract}
Every 2-connected graph outside the three exceptional families satisfies
η(G) ≥ min\{2n, 3n−10\}, proving the conjecture for n ≥ 10. This
candidate proof was generated by GPT-6 Astra in a single model pass for
the Graph Theory AI project. It was selected from the campaign because
the model marked it \texttt{would\_publish:\ true}; it has not been
checked by the adversarial referee, a human mathematician, or a novelty
search.
\end{abstract}

\section{The problem and the claimed
result}\label{the-problem-and-the-claimed-result}

This document concerns
\href{https://graph-theory-ai.github.io/graph-conjectures/arxiv/1602.05184__00/}{Conjecture
5} from \emph{On the difference between the Szeged and Wiener index}.

\subsection{Problem statement}\label{problem-statement}

Conjecture. Let \(G\) be a 2-connected graph of order \(n \geq 10\) not
isomorphic to \(K_n\), \(K^2_n\), or \(K^{n-2}_n\). Then
\(\eta(G) \geq 2n\).

\subsection{Catalog context}\label{catalog-context}

The authors note that while \(\eta(K_n)=0\) and
\(\eta(K^2_n)=\eta(K^{n-2}_n)=2n-6\), computer search found only
finitely many 2-connected graphs with \(\eta(G)<2n\) and none on 10 or
more vertices, motivating this stronger bound beyond the three
exceptional families.

\subsection{Claimed result}\label{claimed-result}

\textbf{Model verdict: proved} (confidence: high).

Every 2-connected graph outside the three exceptional families satisfies
η(G) ≥ min\{2n, 3n−10\}, proving the conjecture for n ≥ 10.

\emph{The model's own caveat:} The proof is self-contained; novelty and
current literature status have not been independently checked.

\section{Model writeup}\label{model-writeup}

The text below is the model output verbatim, apart from moving its
machine-readable verdict block into the summary above and shifting
Markdown heading levels for this document. Mechanical TeX defects and
equation tags were normalized where needed for compilation.

\subsection{1. Statement}\label{statement}

All graphs below are finite and simple. Write \(K_n^t\) for a clique of
order \(n-1\) together with an additional vertex adjacent to exactly
\(t\) clique vertices.

I prove the following stronger statement.

\subsubsection{Theorem}\label{theorem}

Let \(G\) be a 2-connected graph of order \(n\ge4\), with \[
G\not\cong K_n,\ K_n^2,\ K_n^{n-2}.
\] Then \[
\boxed{\eta(G)\ge \min\{2n,\,3n-10\}.}
\qquad\text{(1)}
\]

Consequently, if \(n\ge10\), then \(\eta(G)\ge2n\).

The universal-vertex argument in the supplied attempt is correct; a
verification is included below. The additional ingredient is a reduction
for graphs without universal vertices. It is important that this
reduction produces a dominated vertex \textbf{whose deletion remains
2-connected}.

\subsection{2. Good edges and dominated-vertex
deletion}\label{good-edges-and-dominated-vertex-deletion}

For an edge \(uv\), let \(n_{uv}(u)\) denote the number of vertices
closer to \(u\) than to \(v\). Thus \[
Sz(G)=\sum_{uv\in E(G)}n_{uv}(u)n_{uv}(v),
\qquad
W(G)=\sum_{\{a,b\}\subseteq V(G)}d(a,b).
\]

An edge \(uv\) is \textbf{good} for \(\{a,b\}\) if, after possibly
exchanging \(u,v\), \[
d(a,u)<d(a,v),
\qquad
d(b,v)<d(b,u).
\] Let \(g_G(a,b)\) be the number of good edges, and put \[
D_G(a,b)=g_G(a,b)-d_G(a,b),
\qquad
c_G(a)=\sum_{b\ne a}D_G(a,b).
\]

Every edge of every shortest \(a\)-\(b\) path is good for \(\{a,b\}\).
In particular, \[
D_G(a,b)\ge0.
\] Double counting gives \[
\boxed{\eta(G)=\sum_{\{a,b\}}D_G(a,b)
              =\frac12\sum_a c_G(a).}
\qquad\text{(2)}
\]

We shall repeatedly use the following elementary deletion facts.

\subsubsection{Lemma 2.1}\label{lemma-2.1}

Suppose that \(G\) is 2-connected, \(|V(G)|\ge4\), and \[
N_G[u]\subseteq N_G[v]
\qquad (u\ne v).
\] Put \(H=G-u\).

\begin{enumerate}
\def\labelenumi{\arabic{enumi}.}
\tightlist
\item
  \(H\) is isometric in \(G\).
\item
  The only possible cut vertex of \(H\) is \(v\).
\item
  If \(q_G(u)\) is the number of incidences \[
  \bigl(\{a,b\},e\bigr)
  \] where \(a,b\in V(H)\) and \(e\) is an edge incident with \(u\) that
  is good for \(\{a,b\}\), then \[
  \eta(G)-\eta(H)=c_G(u)+q_G(u).
  \qquad\text{(3)}
  \]
\item
  If \(u\) is simplicial, then \(q_G(u)=0\). If \(u\) is not simplicial,
  then \(q_G(u)\ge2\).
\end{enumerate}

\paragraph{Proof}\label{proof}

Any portion \(x-u-y\) of a path between vertices of \(H\) can be
replaced by a walk through \(v\) of length at most two. This proves that
deleting \(u\) preserves all distances between the remaining vertices.

Suppose \(w\ne v\) were a cut vertex of \(H\). All vertices of
\(N_G(u)\setminus\{w\}\) lie in the component of \(H-w\) containing
\(v\), since they are adjacent to \(v\). Adding \(u\) therefore cannot
connect the components of \(H-w\), contradicting the connectivity of
\(G-w\). This proves the second assertion.

Because distances in \(H\) are unchanged, an edge of \(H\) is good for
an old pair in \(G\) exactly when it is good for that pair in \(H\). The
additional contributions are precisely the pairs containing \(u\), and
the incidences counted by \(q_G(u)\). Hence (3).

If \(u\) is simplicial, then for every \(x\in N(u)\) and every
\(a\ne u\), \[
d(a,x)\le d(a,u).
\] Indeed, take the last neighbor of \(u\) on a shortest \(a\)-\(u\)
path and use that \(N(u)\) is a clique. Thus no vertex other than \(u\)
is closer to \(u\) than to \(x\), and \(ux\) cannot be good for a pair
avoiding \(u\).

Finally, if \(u\) is not simplicial, choose nonadjacent \(x,y\in N(u)\).
Both \(xu\) and \(uy\) are good for \(\{x,y\}\), giving \(q_G(u)\ge2\).
\(\square\)

We also use the standard consequence that deleting a simplicial vertex
from a 2-connected graph of order at least four preserves
2-connectivity. For completeness, if deletion created a cut vertex
\(w\), all remaining neighbors of the simplicial vertex would lie in one
component after deleting \(w\), so restoring it could not reconnect the
graph.

\subsection{3. The structure of a vertex with contribution at most
three}\label{the-structure-of-a-vertex-with-contribution-at-most-three}

The next lemma is the main structural step.

\subsubsection{Lemma 3.1 --- Low-contribution
structure}\label{lemma-3.1-low-contribution-structure}

Let \(G\) be 2-connected, and let \(a\) be a non-universal vertex
satisfying \(c_G(a)\le3\). Put \(A=N_G(a)\).

Every vertex has distance at most two from \(a\), and one of the
following descriptions holds.

\paragraph{Type I}\label{type-i}

There are adjacent vertices \(p,q\in A\) and a vertex \(z\notin N[a]\)
such that \[
V(G)\setminus N[a]=\{z\}\mathbin{\dot\cup}X,
\qquad
N(z)\cap A=\{p,q\}.
\] Either \(X=\varnothing\), or:

\begin{itemize}
\tightlist
\item
  \(G[X]\) is connected;
\item
  every vertex of \(X\) has neighborhood in \(A\) exactly \(\{p\}\);
\item
  there is exactly one edge between \(z\) and \(X\).
\end{itemize}

\paragraph{Type II}\label{type-ii}

There are distinct \(p,q\in A\) and nonempty sets \(X,Y\) such that \[
V(G)\setminus N[a]=X\mathbin{\dot\cup}Y,
\] where:

\begin{itemize}
\tightlist
\item
  \(G[X]\) and \(G[Y]\) are connected;
\item
  every vertex of \(X\) has neighborhood in \(A\) exactly \(\{p\}\);
\item
  every vertex of \(Y\) has neighborhood in \(A\) exactly \(\{q\}\);
\item
  there is exactly one edge \(xy\) between \(X\) and \(Y\).
\end{itemize}

The edges inside \(A\), and inside the indicated connected sets, are
otherwise unrestricted.

Moreover, in any 2-connected graph, \[
a\text{ non-universal}\quad\Longrightarrow\quad c_G(a)\ge2.
\qquad\text{(4)}
\]

\paragraph{Proof}\label{proof-1}

Use distance layers \(L_i\) from \(a\). A neighbor of a vertex in the
preceding layer will be called a \textbf{predecessor}.

Let \(M\) be the set of vertices having more than one shortest path from
\(a\). If \(b\in M\), the union of two distinct shortest \(a\)-\(b\)
paths contains at least \(d(a,b)+2\) edges, all good for \(\{a,b\}\).
Hence \[
D_G(a,b)\ge2.
\qquad\text{(5)}
\] It follows that \(|M|\le1\).

We first record another useful observation. Suppose \(x,y\in L_i\) have
unique shortest paths from \(a\), their predecessors are different, and
\(xy\in E(G)\). If \(p(y)\) denotes the predecessor of \(y\), then \[
d(x,y)=1,\qquad d(x,p(y))=2.
\] Thus \(p(y)y\) is good for \(\{a,x\}\), and it does not belong to the
unique shortest \(a\)-\(x\) path. The analogous assertion holds with
\(x,y\) exchanged.

For a fixed \(x\), different such neighbors \(y\) give different extra
good edges.

\subsubsection{\texorpdfstring{Case 1:
\(M=\varnothing\)}{Case 1: M=\textbackslash varnothing}}\label{case-1-mvarnothing}

The edges between consecutive layers form a rooted tree \(T\): every
non-root vertex has exactly one predecessor.

Let \(F\) be the set of same-layer edges whose endpoints have different
predecessors. The preceding observation gives \[
c_G(a)\ge 2|F|.
\qquad\text{(6)}
\] Thus \(|F|\le1\).

If \(F=\varnothing\), all non-tree edges join siblings. In such a graph,
the proper descendants of any non-root tree vertex have no neighbors
outside that subtree except its root. Consequently, a vertex at distance
at least two from \(a\) would force a cut vertex. Since \(G\) is
2-connected, \(a\) would be universal, contrary to assumption.

Therefore \(F=\{xy\}\), say \(x,y\in L_i\), with distinct predecessors
\(p,q\).

We claim that \(p,q\) have the same predecessor. Otherwise \(p,q\) are
nonadjacent, since \(pq\) would be a second edge of \(F\). They also
have no common neighbor:

\begin{itemize}
\tightlist
\item
  a common neighbor in \(L_{i-2}\) contradicts their distinct
  predecessors;
\item
  a common neighbor in \(L_{i-1}\) would produce another edge of \(F\);
\item
  a common neighbor in \(L_i\) would have two shortest paths from \(a\).
\end{itemize}

Hence \[
d(p,q)=3
\] via \(p-x-y-q\), while \(d(p,y)=d(q,x)=2\). Therefore \(qy\) is an
extra good edge for \(\{a,p\}\), and \(px\) is an extra good edge for
\(\{a,q\}\). Together with the extra contributions at \(x\) and \(y\),
this gives \(c_G(a)\ge4\), a contradiction.

Let \(w\) be the common predecessor of \(p,q\). Apart from sibling
edges, the only non-tree edge is \(xy\), which lies entirely among
proper descendants of \(w\). If \(w\ne a\), those descendants are
separated from \(a\) by \(w\), contradicting 2-connectivity. Thus
\(w=a\), and \(i=2\).

There can be no vertex in \(L_3\) or a later layer. Indeed, the proper
descendants of its ancestor in \(L_2\) would have no exit except through
that ancestor: the sole exceptional edge \(xy\) has both endpoints in
\(L_2\).

Every vertex of \(L_2\) therefore has a unique neighbor in \(A\). The
only edge between different predecessor classes is \(xy\). A class other
than those of \(p,q\), or a component of the \(p\)-class not containing
\(x\), would be separated by its predecessor. The analogous statement
holds for \(q\). This gives Type II.

\subsubsection{\texorpdfstring{Case 2:
\(M=\{z\}\)}{Case 2: M=\textbackslash\{z\textbackslash\}}}\label{case-2-mz}

The vertex \(z\) has at least two predecessors; otherwise its unique
predecessor would also belong to \(M\). It has no neighbor in the next
layer, since such a neighbor would also belong to \(M\).

If \(z\) had \(r\ge3\) predecessors, the last two layers of its shortest
paths would contain at least \(2r\) good edges, only two of which belong
to one fixed shortest path. Hence \[
D_G(a,z)\ge2r-2\ge4,
\] which is impossible. Thus \(z\) has exactly two predecessors, say
\(p,q\).

If \(pq\notin E(G)\), then \(qz\) is an extra good edge for \(\{a,p\}\),
and \(pz\) is one for \(\{a,q\}\). Together with \(D_G(a,z)\ge2\), this
again gives \(c_G(a)\ge4\). Hence \(pq\in E(G)\).

There cannot be a same-layer edge between two uniquely reached vertices
with different predecessors: it would contribute at least one at each
endpoint in addition to the contribution of at least two at \(z\). In
particular, \(p,q\) have the same predecessor, say \(w\).

Consider a same-layer neighbor \(t\) of \(z\). Its unique predecessor
must be \(p\) or \(q\). Otherwise both \(pz\) and \(qz\) are extra good
edges for \(\{a,t\}\), giving \(D_G(a,t)\ge2\). If its predecessor is
\(p\) or \(q\), it still contributes at least one. Thus \(z\) has at
most one same-layer neighbor.

Choose a breadth-first tree by making \(p\) the predecessor of \(z\).
All non-tree edges are sibling edges, except \(qz\) and the possible
same-layer edge incident with \(z\). These exceptional edges lie
entirely among proper descendants of \(w\). As in Case 1, 2-connectivity
forces \(w=a\). Thus \(z\in L_2\). The same subtree-separation argument
rules out all later layers.

If \(z\) has no same-layer neighbor, every other predecessor class in
\(L_2\) would be separated by its predecessor, so \(L_2=\{z\}\).

Otherwise let \(t\) be its sole same-layer neighbor, and exchange
\(p,q\) if necessary so that \(p\) is the predecessor of \(t\). All
remaining vertices of \(L_2\) must have predecessor \(p\), and their
induced graph must be connected and contain \(t\); any other component
would be separated by its predecessor. This is Type I.

Finally, if \(c_G(a)\le1\), then (5) forces \(M=\varnothing\), and (6)
forces \(F=\varnothing\). The argument in Case 1 then makes \(a\)
universal. This proves (4). \(\square\)

\subsection{4. Consequences of the low-contribution
structure}\label{consequences-of-the-low-contribution-structure}

\subsubsection{Lemma 4.1 --- A 2-connected deletion
exists}\label{lemma-4.1-a-2-connected-deletion-exists}

Let \(G\) be 2-connected with no universal vertex. If some vertex has
contribution at most three, then either \(G\cong C_5\), or there are
distinct vertices \(u,v\) such that \[
N[u]\subseteq N[v]
\quad\text{and}\quad
G-u\text{ is 2-connected}.
\qquad\text{(7)}
\]

\paragraph{Proof}\label{proof-2}

Apply Lemma 3.1.

\paragraph{Type I}\label{type-i-1}

If \(X=\varnothing\), then \(z\) is simplicial, with adjacent neighbors
\(p,q\), so deleting \(z\) preserves 2-connectivity.

Suppose \(X\ne\varnothing\), and let \(zt\) be its sole edge to \(z\).
Every vertex of \(X\) is dominated by \(p\).

If \(X=\{t\}\), then \(t\) is simplicial, with adjacent neighbors
\(p,z\), and can be deleted.

Otherwise choose \(u\in X\setminus\{t\}\) such that \(G[X]-u\) remains
connected. Such a vertex exists by taking a leaf other than \(t\) in a
spanning tree of \(G[X]\).

The graph \(G-p\) is connected, and \(X\) is attached to its complement
there only through \(tz\). Deleting \(u\) preserves this connectivity.
Thus \(G-\{u,p\}\) is connected. Since \(u\) is dominated by \(p\),
Lemma 2.1 shows that \(G-u\) is 2-connected.

\paragraph{Type II}\label{type-ii-1}

If \(|X|\ge2\), choose a non-cut vertex \(u\ne x\) of \(G[X]\). It is
dominated by \(p\). In \(G-p\), the set \(X\) is attached to its
complement only by \(xy\), and deleting \(u\) preserves connectivity.
Lemma 2.1 again gives a 2-connected deletion. The case \(|Y|\ge2\) is
symmetric.

We may therefore assume \(X=\{x\}\) and \(Y=\{y\}\).

If \(A=\{p,q\}\), then:

\begin{itemize}
\tightlist
\item
  if \(pq\notin E(G)\), the graph is \(C_5\);
\item
  if \(pq\in E(G)\), the vertex \(a\) is simplicial and can be deleted.
\end{itemize}

Finally suppose \(|A|\ge3\). The graph \[
R=G[A]+pq
\] is connected, because \(G-a\) is connected and the path \(p-x-y-q\)
connects \(p\) to \(q\).

Choose a spanning tree of \(R\) containing \(pq\). Since this tree has
at least three vertices, it has a leaf \(u\notin\{p,q\}\). Thus \(R-u\)
is connected. The vertex \(u\) is dominated by \(a\), and replacing the
added edge \(pq\) by \(p-x-y-q\) shows that \(G-\{a,u\}\) is connected.
Lemma 2.1 proves that \(G-u\) is 2-connected. \(\square\)

The next lemma prevents low-increment simplicial deletions in a
counterexample without a universal vertex.

\subsubsection{Lemma 4.2 --- Low-contribution simplicial
vertices}\label{lemma-4.2-low-contribution-simplicial-vertices}

Let \(G\) be 2-connected and have no universal vertex. If \(a\) is
simplicial and \(c_G(a)\le3\), then \[
\eta(G)\ge2|V(G)|.
\qquad\text{(8)}
\]

\paragraph{Proof}\label{proof-3}

In Lemma 3.1, \(A=N(a)\) is now a clique.

Type I is impossible: the vertex \(p\) is adjacent to \(a\), all of
\(A\), \(z\), and all of \(X\), so it is universal.

Thus Type II holds. Put \[
K=\{a\}\cup A.
\] The vertices \(p,x,y,q\) induce the cycle \[
p-x-y-q-p.
\] It is isometric, so the sum of \(D_G\) over pairs within these four
vertices is at least \[
\eta(C_4)=8.
\qquad\text{(9)}
\]

We show that every other vertex supplies at least two additional units,
using only pairs with these four vertices.

\paragraph{\texorpdfstring{Vertices in
\(K\setminus\{p,q\}\)}{Vertices in K\textbackslash setminus\textbackslash\{p,q\textbackslash\}}}\label{vertices-in-ksetminuspq}

For \(b\in K\setminus\{p,q\}\), the path \(b-p-x\) is shortest, and
\(qy\) is an additional good edge for \(\{b,x\}\). Hence \[
D_G(b,x)\ge1.
\] Symmetrically, \[
D_G(b,y)\ge1.
\]

\paragraph{\texorpdfstring{Vertices in
\(X\setminus\{x\}\)}{Vertices in X\textbackslash setminus\textbackslash\{x\textbackslash\}}}\label{vertices-in-xsetminusx}

Let \(b\in X\setminus\{x\}\).

If \(bx\in E(G)\), then:

\begin{itemize}
\tightlist
\item
  \(b-p-q\) is shortest and \(xy\) is an additional good edge for
  \(\{b,q\}\);
\item
  \(b-x-y\) is shortest and \(pq\) is an additional good edge for
  \(\{b,y\}\).
\end{itemize}

Thus \(D_G(b,q)+D_G(b,y)\ge2\).

If \(bx\notin E(G)\), then \(d(b,y)=3\), and both \[
b-p-q-y,\qquad b-p-x-y
\] are shortest paths. Their union has five edges, so \[
D_G(b,y)\ge2.
\]

The argument for \(Y\setminus\{y\}\) is symmetric.

These counted pairs are distinct: each has one endpoint outside the
four-cycle and one inside it. All uncounted contributions are
nonnegative. Therefore \[
\eta(G)\ge8+2(|V(G)|-4)=2|V(G)|.
\] \(\square\)

Combining the preceding lemmas gives the reduction needed for induction.

\subsubsection{Proposition 4.3 --- Reduction with increment at least
four}\label{proposition-4.3-reduction-with-increment-at-least-four}

Let \(G\) be 2-connected, with no universal vertex, and suppose \[
\eta(G)<2|V(G)|.
\] Then either \(G\cong C_5\), or there is a vertex \(u\) dominated by
another vertex such that \(H=G-u\) is 2-connected and \[
\boxed{\eta(G)-\eta(H)\ge4.}
\qquad\text{(10)}
\]

\paragraph{Proof}\label{proof-4}

By (2), some vertex has contribution at most three. Unless \(G=C_5\),
Lemma 4.1 supplies the required dominated vertex \(u\) with a
2-connected deletion.

If \(u\) is not simplicial, then Lemma 3.1 gives \(c_G(u)\ge2\), and
Lemma 2.1 gives \(q_G(u)\ge2\). Hence the increment is at least four.

If \(u\) is simplicial, Lemma 4.2 and the hypothesis \(\eta(G)<2|V(G)|\)
force \(c_G(u)\ge4\). Equation (3) again gives the result. \(\square\)

\subsection{5. Graphs with a universal
vertex}\label{graphs-with-a-universal-vertex}

This section verifies the relevant result from the supplied attempt.

For a graph \(Q\), let \(p_3(Q),p_4(Q),c_4(Q)\) count induced copies of
\(P_3,P_4,C_4\), respectively.

\subsubsection{Lemma 5.1}\label{lemma-5.1}

If \(\operatorname{diam}(G)\le2\), then \[
\eta(G)
=
2\bigl(p_3(G)-\overline m(G)\bigr)
+p_4(G)+4c_4(G),
\qquad\text{(11)}
\] where \(\overline m(G)\) is the number of nonedges.

\paragraph{Proof}\label{proof-5}

For a nonedge \(ab\), every common neighbor gives two good edges. Every
other good edge is the middle edge of an induced \(P_4\) with endpoints
\(a,b\). Thus \[
D_G(a,b)
=
2\bigl(|N(a)\cap N(b)|-1\bigr)
+\#\{\text{induced \(P_4\)'s with endpoints \(a,b\)}\}.
\]

For an edge \(ab\), every good edge other than \(ab\) is the edge
opposite \(ab\) in an induced \(C_4\). Summing over pairs counts each
induced \(C_4\) four times. Summing common-neighbor counts over nonedges
counts \(p_3(G)\). This proves (11). \(\square\)

In particular, for a cone \(G=K_1\vee Q\), \[
\boxed{\eta(G)=2p_3(Q)+p_4(Q)+4c_4(Q).}
\qquad\text{(12)}
\] Indeed, the apex creates one additional induced \(P_3\) for each
nonedge of \(Q\), and belongs to no induced \(P_4\) or \(C_4\).

\subsubsection{Lemma 5.2}\label{lemma-5.2}

Let \(Q\) be a connected, noncomplete graph of order \(m\). Unless \(Q\)
is \(K_m-e\) or a clique \(K_{m-1}\) with one pendant vertex, \[
2p_3(Q)+p_4(Q)+4c_4(Q)\ge3m-7.
\qquad\text{(13)}
\] For each of the two exceptional graphs, the left side equals
\(2m-4\).

\paragraph{Proof}\label{proof-6}

Write \[
F(Q)=2p_3(Q)+p_4(Q)+4c_4(Q).
\] Choose a maximum clique \(C\), of order \(c\ge2\). Let \(L_i\) now
denote distance layers from \(C\), and put \[
\ell=|L_1|,
\qquad
r=\sum_{i\ge2}|L_i|.
\] Thus \(m=c+\ell+r\), and \(\ell\ge1\).

For \(x\in L_1\), set \(a_x=|N(x)\cap C|\). Maximality of \(C\) gives \[
1\le a_x\le c-1.
\] The vertex \(x\), one neighbor in \(C\), and one nonneighbor in \(C\)
give \[
a_x(c-a_x)\ge c-1
\] induced \(P_3\)'s.

Every vertex in \(L_i\), \(i\ge2\), supplies another induced \(P_3\),
using the first three vertices of a shortest path to \(C\). These copies
are distinct, by their unique vertex of maximum layer. Consequently, \[
p_3(Q)\ge(c-1)\ell+r.
\qquad\text{(14)}
\]

Likewise, \[
p_4(Q)\ge r.
\qquad\text{(15)}
\] For a vertex at distance at least three, take four vertices of a
shortest path. For \(x\in L_2\), take \(x-y-v\) with \(v\in C\), and
choose \(w\in C\setminus N(y)\); then \(x-y-v-w\) is induced. Again the
maximum-layer vertex distinguishes the copies.

\paragraph{\texorpdfstring{Case 1: \(c\ge3\) and
\(\ell\ge2\)}{Case 1: c\textbackslash ge3 and \textbackslash ell\textbackslash ge2}}\label{case-1-cge3-and-ellge2}

Equations (14)--(15) give \[
F(Q)\ge2(c-1)\ell+3r.
\] The difference from \(3m-7\) is \[
(2c-5)\ell-3c+7\ge c-3\ge0.
\]

\paragraph{\texorpdfstring{Case 2: \(c\ge3\) and
\(\ell=1\)}{Case 2: c\textbackslash ge3 and \textbackslash ell=1}}\label{case-2-cge3-and-ell1}

Write \(L_1=\{x\}\), \(a=|N(x)\cap C|\), and \(A=a(c-a)\ge c-1\).

If \(r=0\), then \(F(Q)=2A\). The cases \(a=1\) and \(a=c-1\) are the
two exceptions. Otherwise \(2\le a\le c-2\), and \[
F(Q)\ge4(c-2)\ge3c-4=3m-7.
\]

Suppose \(r>0\), and write \[
s=|L_2|,\qquad t=\sum_{i\ge3}|L_i|.
\] Every vertex of \(L_2\) is adjacent to \(x\). It supplies \(A\)
induced paths \(y-x-v-w\), where \(v\in N(x)\cap C\) and
\(w\in C\setminus N(x)\). Hence \[
p_3(Q)\ge A+s+t,
\qquad
p_4(Q)\ge As+t.
\] Therefore \[
F(Q)\ge2A+(A+2)s+3t.
\] Subtracting \(3m-7=3c+3s+3t-4\), and using \(s\ge1\), gives \[
2A+(A-1)s-3c+4
\ge2(c-1)+(c-2)-3c+4=0.
\]

\paragraph{\texorpdfstring{Case 3:
\(c=2\)}{Case 3: c=2}}\label{case-3-c2}

Write \(C=\{u,v\}\). The graph is triangle-free. Partition \(L_1\) into
\[
A=N(u)\cap L_1,\qquad B=N(v)\cap L_1,
\] with sizes \(\alpha,\beta\), so \(\ell=\alpha+\beta\). Let \(e\) be
the number of edges between \(A\) and \(B\).

Besides the copies in (14), pairs in \(A\) and in \(B\) give \[
\binom{\alpha}{2}+\binom{\beta}{2}
\] additional induced \(P_3\)'s. A nonedge between \(A,B\) gives an
induced \(P_4\) through \(u,v\), and an edge between them gives an
induced \(C_4\). These copies are disjoint from those counted in (15).
Thus \begin{equation*}
\begin{aligned}
F(Q)
&\ge
2\left(\ell+r+\binom{\alpha}{2}+\binom{\beta}{2}\right)
+r+(\alpha\beta-e)+4e\\
&=
2\ell+3r+
\alpha(\alpha-1)+\beta(\beta-1)+\alpha\beta+3e.
\end{aligned}
\end{equation*} For nonnegative integers \(\alpha,\beta\) with
\(\alpha+\beta\ge1\), \[
\alpha(\alpha-1)+\beta(\beta-1)+\alpha\beta
\ge\alpha+\beta-1.
\] If both are positive, this follows by bounding the first two terms by
\(\alpha-1,\beta-1\), and \(\alpha\beta\) by \(1\); if one is zero it is
immediate. Hence \[
F(Q)\ge3\ell+3r-1=3m-7.
\]

Finally, each exceptional graph has \(m-2\) induced \(P_3\)'s and no
induced \(P_4\) or \(C_4\), giving \(F(Q)=2m-4\). \(\square\)

\subsubsection{Proposition 5.3}\label{proposition-5.3}

Let \(G\) be 2-connected of order \(n\), with a universal vertex. If \[
G\not\cong K_n,\ K_n^2,\ K_n^{n-2},
\] then \[
\boxed{\eta(G)\ge3n-10.}
\qquad\text{(16)}
\]

\paragraph{Proof}\label{proof-7}

Write \(G=K_1\vee Q\). The graph \(Q\) is connected, since deleting the
universal vertex preserves connectivity.

If \(Q\) is complete, \(G=K_n\). The two exceptions of Lemma 5.2 give,
respectively, \(K_n^{n-2}\) and \(K_n^2\). Otherwise (12)--(13), with
\(m=n-1\), give \[
\eta(G)\ge3(n-1)-7=3n-10.
\] \(\square\)

\subsection{6. Proof of the theorem}\label{proof-of-the-theorem}

First, formula (11) gives \[
\boxed{\eta(K_m^t)=2(m-1-t)(t-1).}
\qquad\text{(17)}
\] Indeed, there are \(m-1-t\) nonedges, each with \(t\) common
neighbors, and no induced \(P_4\) or \(C_4\). In particular, \[
\eta(K_m^2)=\eta(K_m^{m-2})=2m-6.
\qquad\text{(18)}
\]

Put \[
f(n)=\min\{2n,3n-10\}.
\] We prove (1) by induction on \(n\).

For \(n=4\), the only admissible graph is \(C_4\), and
\(\eta(C_4)=8\ge f(4)\).

Now let \(n\ge5\), and suppose the assertion holds at smaller orders.

If \(G\) has a universal vertex, Proposition 5.3 gives the result. We
may therefore assume that \(G\) has no universal vertex.

If \(\eta(G)\ge2n\), there is nothing to prove. If \(G=C_5\), then \[
\eta(G)=20-15=5=f(5).
\] Otherwise Proposition 4.3 supplies a dominated vertex \(u\) such that
\[
H=G-u
\] is 2-connected and \[
\eta(G)\ge\eta(H)+4.
\qquad\text{(19)}
\] Write \(m=n-1\).

\subsubsection{\texorpdfstring{Case 1: \(H\) is not
exceptional}{Case 1: H is not exceptional}}\label{case-1-h-is-not-exceptional}

By induction, \[
\eta(G)\ge f(n-1)+4.
\] But \[
f(n-1)+4
=\min\{2n+2,3n-9\}
\ge\min\{2n,3n-10\}=f(n).
\] So this case is finished.

It remains to check every possible exceptional predecessor.

\subsubsection{\texorpdfstring{Case 2:
\(H\cong K_m\)}{Case 2: H\textbackslash cong K\_m}}\label{case-2-hcong-k_m}

Every neighbor of \(u\) is universal in \(G\), contrary to our
assumption that \(G\) has no universal vertex.

\subsubsection{\texorpdfstring{Case 3:
\(H\cong K_m^{m-2}=K_m-e\)}{Case 3: H\textbackslash cong K\_m\^{}\{m-2\}=K\_m-e}}\label{case-3-hcong-k_mm-2k_m-e}

Let \(x,y\) be the endpoints of its missing edge. Every vertex of
\(H\setminus\{x,y\}\) is universal in \(H\). Since \(G\) has no
universal vertex, none is adjacent to \(u\).

Because \(G\) is 2-connected, \(d_G(u)\ge2\), so \[
N_G(u)=\{x,y\}.
\] But neither \(x\) nor \(y\) can dominate \(u\), since
\(xy\notin E(G)\). This contradicts the choice of \(u\).

\subsubsection{\texorpdfstring{Case 4:
\(H\cong K_m^2\)}{Case 4: H\textbackslash cong K\_m\^{}2}}\label{case-4-hcong-k_m2}

The case \(m=4\) is already included in Case 3. Otherwise write \(H\) as
a clique \(B\) of order \(m-1\) and an additional vertex \(z\) with \[
N_H(z)=\{p,q\}\subseteq B.
\] Put \[
R=B\setminus\{p,q\}.
\]

The vertices \(p,q\) are universal in \(H\). Thus \(u\) is adjacent to
neither of them.

We claim that \(u\) is not adjacent to \(z\). Otherwise, since
\(d_G(u)\ge2\), it also has a neighbor in \(R\). A vertex dominating
\(u\) would have to be a neighbor of \(u\); but \(z\) is adjacent to no
vertex of \(R\), and no vertex of \(R\) is adjacent to \(z\). No
possible dominator exists.

Consequently, \[
N_G(u)\subseteq R.
\] Let \(t=|N_G(u)|\ge2\). The vertex \(u\) is simplicial.

Each of \(p,q\) has all \(t\) vertices of \(N(u)\) as common neighbors
with \(u\), so \[
D_G(u,p)\ge2(t-1),
\qquad
D_G(u,q)\ge2(t-1).
\qquad\text{(20)}
\]

Also \(d_G(u,z)=3\). For every \(r\in N(u)\), the two paths \[
u-r-p-z,\qquad u-r-q-z
\] are shortest. Their union over all \(r\) has \(3t+2\) edges, all good
for \(\{u,z\}\). Therefore \[
D_G(u,z)\ge3t-1.
\qquad\text{(21)}
\] Combining (20)--(21), \[
c_G(u)\ge7t-5\ge9.
\] By Lemma 2.1 and (18), \[
\eta(G)\ge(2m-6)+9=2n+1.
\] This contradicts the case assumption \(\eta(G)<2n\).

All exceptional predecessors have now been covered, completing the
induction and proving (1). \(\square\)

Since \(3n-10\ge2n\) exactly when \(n\ge10\), the conjecture follows.

\subsection{7. Sharpness and scope}\label{sharpness-and-scope}

The stronger bound is sharp for every \(n\ge5\).

\paragraph{\texorpdfstring{The \(3n-10\)
branch}{The 3n-10 branch}}\label{the-3n-10-branch}

For \[
G=K_1\vee P_{n-1},
\] equation (12) gives \[
\eta(G)=2(n-3)+(n-4)=3n-10.
\] These graphs are admissible for \(n\ge5\).

\paragraph{\texorpdfstring{The \(2n\)
branch}{The 2n branch}}\label{the-2n-branch}

Start with the cycle \(p-x-y-q-p\), and enlarge its edge \(pq\) to a
clique by adding \(n-4\) vertices, each adjacent to every vertex already
in that clique and to neither \(x\) nor \(y\).

Equivalently, take a clique of order \(n-2\), an additional edge \(xy\),
and exactly two cross edges forming a matching.

Each added vertex is simplicial and has contribution exactly two: its
only nonzero contributions are one with \(x\) and one with \(y\). Lemma
2.1 therefore gives \[
\eta(G)=\eta(C_4)+2(n-4)=2n.
\]

Thus the transition at \(n=10\) is genuine.

No computational enumeration or unproved external assertion is used. In
particular, the previously untreated graphs without universal vertices
are covered by Proposition 4.3, and the induction explicitly handles all
three exceptional predecessor families.

\appendix

\section{Provenance and review
status}\label{provenance-and-review-status}

\begin{warning}

\textbf{UNREFEREED MODEL OUTPUT.} This document was selected solely
because GPT-6 Astra labelled its own result
\texttt{would\_publish:\ true} and returned \texttt{proved} or
\texttt{disproved}. It has not passed the adversarial LLM referee used
for the earlier Sol campaign, has not been checked by a human
mathematician, and has not been checked for novelty. Treat every
mathematical and bibliographic claim as unverified.

\end{warning}

{\def\LTcaptype{none} % do not increment counter
\begin{longtable}[]{@{}
  >{\raggedright\arraybackslash}p{(\linewidth - 2\tabcolsep) * \real{0.5000}}
  >{\raggedright\arraybackslash}p{(\linewidth - 2\tabcolsep) * \real{0.5000}}@{}}
\toprule\noalign{}
\endhead
\bottomrule\noalign{}
\endlastfoot
Catalog id & \texttt{1602.05184\_\_00} \\
Catalog entry &
\href{https://graph-theory-ai.github.io/graph-conjectures/arxiv/1602.05184__00/}{Conjecture
5} \\
Source corpus & arxiv \\
Campaign leg & selected second attempts (\texttt{attacks\_retry}) \\
Source paper / entry & On the difference between the Szeged and Wiener
index \\
Model verdict & \textbf{proved} (confidence: high) \\
Selection criterion & \texttt{would\_publish:\ true} and verdict
\texttt{proved} or \texttt{disproved} \\
Model & \texttt{gpt-6-astra}, reasoning effort \texttt{max},
\texttt{mode=pro}, flex service tier \\
Original artifact &
\texttt{attacks\_retry/1602.05184\_\_00/output.md} \\
Independent review & \textbf{None} \\
arXiv & \href{https://arxiv.org/abs/1602.05184}{arXiv:1602.05184} \\
Previous attempt & \texttt{attacks/1602.05184\_\_00}: gpt-5.6-sol
verdict \texttt{partial} \\
Model's caveats & The proof is self-contained; novelty and current
literature status have not been independently checked. \\
\end{longtable}
}

The generated Markdown and LaTeX sources for this PDF are stored under
\texttt{to\_review\_astra/src/1602.05184\_\_00/}.

\end{document}
